\thispagestyle{pokedex}
\pdfbookmark[0]{Control del documento}{control}

\vspace*{6pt}
{\sffamily\bfseries\LARGE\color{pkInk} Control del documento\par}
\vspace{4pt}{\color{pkRed}\rule{2.2cm}{2pt}}\par
\vspace{18pt}

\par\addvspace{8pt}\noindent\begin{pkbox}{colspec={X[2,l]X[5,l]}}
  Campo & Valor \\
  Título & Manual Técnico — PokéDex Manager \\
  Versión & 1.0 \\
  Fecha & 13 de agosto de 2026 \\
  Estado & Vigente \\
  Clasificación & Uso interno \\
  Repositorio & \pkpath{https://github.com/Skulltulaidm/pokedex-manager} (público) \\
  Rama documentada & \pkpath{main} \\
  Revisión del código & \pkpath{c5d3bd3} \\
  Audiencia & Ingeniería de software, QA, operaciones \\
  Despliegue en producción & \pkpath{https://web-production-78a4b.up.railway.app} \\
  API en producción & \pkpath{https://api-production-81ae.up.railway.app} \\
  Documento hermano & Manual de usuario (mismo repositorio) \\
\end{pkbox}\par\addvspace{8pt}

\vspace{20pt}
{\pksb\large\color{pkInk} Historial de revisiones\par}
\vspace{8pt}

\par\addvspace{8pt}\noindent\begin{pkbox}{colspec={X[1,l]X[1.4,l]X[4,l]X[1.4,l]}}
  Versión & Fecha & Cambio & Estado \\
  1.0 & 2026-08-13 & Primera emisión. Cubre arquitectura, modelo de datos, backend, frontend, requisitos, QA, CI/CD, instalación y bitácora de decisiones. & Vigente \\
\end{pkbox}\par\addvspace{8pt}

\vspace{20pt}
{\pksb\large\color{pkInk} Qué contiene y qué no\par}
\vspace{8pt}

\begin{pknote}[Alcance]
Este documento describe \textbf{el sistema tal como está construido}, verificado
leyendo el código de la revisión indicada arriba y ejecutando sus suites de
pruebas. Las cifras que aparecen (número de rutas, de tablas, de pruebas, de
herramientas MCP) se obtuvieron de esa revisión, no de una estimación.

Fuera de alcance: la operación día a día del producto para una persona usuaria
—eso está en el manual de usuario— y cualquier funcionalidad planificada que
todavía no exista en el repositorio.
\end{pknote}

\vspace{6pt}
{\pksb\large\color{pkInk} Convenciones\par}
\vspace{8pt}

\begin{itemize}[leftmargin=1.1em,itemsep=2pt,topsep=2pt]
  \item \pkpath{ruta/al/archivo.py} es una ruta relativa a la raíz del repositorio.
  \item \pkcode{identificador} es una función, clase, columna o variable de entorno tal cual aparece en el código.
  \item Las citas de código son literales; los comentarios que las acompañan también están en el repositorio.
  \item Los identificadores de requisito siguen \pkid{RF-nn} (funcional), \pkid{RNF-nn} (no funcional) y \pkid{HU-nn} (historia de usuario).
\end{itemize}
